\chapter{Definite Integrals and the Fundamental Theorem}

The definite integral $\int_a^b f(x)\,\mathrm{d}x$ is the signed area
between the graph of $f$ and the $x$-axis on $[a,b]$. The Fundamental
Theorem of Calculus is the bridge between this
geometric object and the indefinite integral of Section 1: and the
reason any of the antiderivatives we have been computing matter for
economic measurement.

\section{Core Rules}

\begin{rulebox}{Fundamental Theorem of Calculus, Part II (evaluation)}
If $f$ is continuous on $[a,b]$ and $F$ is any antiderivative of $f$ on
$[a,b]$, then
\[
  \int_{a}^{b} f(x)\,\mathrm{d}x \;=\; F(b)-F(a) \;=:\; \bigl[F(x)\bigr]_{a}^{b} .
\]
\end{rulebox}

\begin{rulebox}{Fundamental Theorem, Part I (differentiation)}
If $f$ is continuous and $A(x):=\int_{a}^{x} f(t)\,\mathrm{d}t$, then
$A'(x)=f(x)$. With a variable upper limit $g(x)$,
\[
  \frac{\mathrm{d}}{\mathrm{d}x}\int_{a}^{g(x)} f(t)\,\mathrm{d}t
    \;=\; f(g(x))\cdot g'(x) .
\]
\end{rulebox}

\begin{rulebox}{Properties of the definite integral}
For continuous $f,g$ on the relevant intervals and constants $\alpha,\beta$:
\begin{itemize}[leftmargin=1.6em,itemsep=0.3ex]
  \item \textbf{Linearity:}
        $\displaystyle\int_{a}^{b}(\alpha f+\beta g)\,\mathrm{d}x
          =\alpha\int_{a}^{b} f\,\mathrm{d}x+\beta\int_{a}^{b} g\,\mathrm{d}x$;
  \item \textbf{Additivity over intervals:} for $a\le c\le b$,
        $\displaystyle\int_{a}^{b} f\,\mathrm{d}x
          =\int_{a}^{c} f\,\mathrm{d}x+\int_{c}^{b} f\,\mathrm{d}x$;
  \item \textbf{Reversal of limits:}
        $\displaystyle\int_{b}^{a} f\,\mathrm{d}x=-\int_{a}^{b} f\,\mathrm{d}x$;
  \item \textbf{Empty interval:}
        $\displaystyle\int_{a}^{a} f\,\mathrm{d}x=0$;
  \item \textbf{Monotonicity:} if $f\le g$ on $[a,b]$, then
        $\displaystyle\int_{a}^{b} f\,\mathrm{d}x\le\int_{a}^{b} g\,\mathrm{d}x$.
\end{itemize}
\end{rulebox}

\begin{rulebox}{Substitution with limits}
If $u=g(x)$ then
\[
  \int_{a}^{b} f(g(x))\,g'(x)\,\mathrm{d}x \;=\; \int_{g(a)}^{g(b)} f(u)\,\mathrm{d}u .
\]
Either change the limits to $u$-values (faster) or substitute back to $x$
before evaluating.
\end{rulebox}

\begin{rulebox}{Average value of a function}
The average (mean) value of $f$ over $[a,b]$ is
\[
  \bar f \;=\; \frac{1}{b-a}\int_{a}^{b} f(x)\,\mathrm{d}x .
\]
In economics this is how one averages a continuously varying cost,
price, or income over time.
\end{rulebox}

\section{Worked Example}

\begin{worked}{Evaluate $\displaystyle\int_{1}^{4}\Bigl(2x-\frac{3}{\sqrt x}\Bigr)\mathrm{d}x$ and its average value; differentiate $\displaystyle\int_{0}^{x^{2}}\mathrm{e}^{-t^{2}}\mathrm{d}t$; and expose the trap in $\displaystyle\int_{-1}^{1}\frac{\mathrm{d}x}{x^{2}}$}
\wstep{Step 1. Antidifferentiate, then evaluate at the two limits.}
The Fundamental Theorem says $\int_{a}^{b}f=F(b)-F(a)$ for any
antiderivative $F$. Rewrite first: $3/\sqrt x=3x^{-1/2}$, so
\[
  F(x)=x^{2}-3\cdot\frac{x^{1/2}}{1/2}=x^{2}-6\sqrt x .
\]
\[
  \int_{1}^{4}\Bigl(2x-\frac{3}{\sqrt x}\Bigr)\mathrm{d}x
  =\bigl[x^{2}-6\sqrt x\bigr]_{1}^{4}
  =(16-12)-(1-6)=4-(-5)=\boxed{9}
\]
\wtrap{$F(b)-F(a)$, and the subtraction is of the \emph{whole} expression.
$(1-6)=-5$, and subtracting $-5$ adds $5$. Bracket the lower value before
subtracting; a sign lost here is the most common definite-integral error.
And no $+C$ is needed --- it cancels in the difference, which is why any
antiderivative works.}

\wstep{Step 2. Average value: divide by the length of the interval.}
\[
  \bar f=\frac{1}{b-a}\int_{a}^{b}f
  =\frac{9}{4-1}=3 .
\]
So the function averages $3$ over $[1,4]$ --- and since $f(1)=-1$ and
$f(4)=6.5$, a value of $3$ in between is plausible. \checkmark\ Dividing by
the interval length is what turns an area into a level.

\wstep{Step 3. The other half of the theorem: differentiating an integral.}
Part~I says
\[
  \frac{\mathrm{d}}{\mathrm{d}x}\int_{a}^{x}f(t)\,\mathrm{d}t=f(x) .
\]
The upper limit here is $x^{2}$, not $x$, so the chain rule intervenes.
Write $G(u)=\int_{0}^{u}\mathrm{e}^{-t^{2}}\mathrm{d}t$ with $u=x^{2}$:
\[
  \frac{\mathrm{d}}{\mathrm{d}x}G(x^{2})=G'(x^{2})\cdot2x
  =\mathrm{e}^{-(x^{2})^{2}}\cdot2x
  =\boxed{2x\,\mathrm{e}^{-x^{4}}}
\]
\wtrap{Two traps at once. The inner derivative $2x$ is required; and $t$ is
replaced by the \emph{whole} upper limit, giving $\mathrm{e}^{-x^{4}}$, not
$\mathrm{e}^{-x^{2}}$. Note also that $\mathrm{e}^{-t^{2}}$ has no
elementary antiderivative --- yet the derivative of its integral is
immediate. Part~I does not need you to be able to integrate.}

\wstep{Step 4. Now the trap. Evaluate $\int_{-1}^{1}x^{-2}\mathrm{d}x$ mechanically.}
\[
  \int_{-1}^{1}\frac{\mathrm{d}x}{x^{2}}
  \;\overset{?}{=}\;\Bigl[-\frac1x\Bigr]_{-1}^{1}
  =-1-(+1)=-2 .
\]
This is nonsense: the integrand $1/x^{2}$ is \textbf{strictly positive}
everywhere it is defined, so no correct answer can be negative.

\wstep{Step 5. Diagnose it.}
The Fundamental Theorem requires $f$ to be continuous on $[a,b]$ and $F$ to
be an antiderivative \emph{throughout}. Here $1/x^{2}$ has a pole at $x=0$,
which lies inside $[-1,1]$, and $-1/x$ is not an antiderivative across it.
The integral is \textbf{improper}, and splitting it shows it diverges:
\[
  \int_{0}^{1}\frac{\mathrm{d}x}{x^{2}}
  =\lim_{\varepsilon\to0^{+}}\Bigl[-\frac1x\Bigr]_{\varepsilon}^{1}
  =\lim_{\varepsilon\to0^{+}}\Bigl(-1+\frac{1}{\varepsilon}\Bigr)=+\infty .
\]
\[
  \boxed{\int_{-1}^{1}\frac{\mathrm{d}x}{x^{2}}\ \text{diverges.}}
\]
\wtrap{Scan the interval for poles of the integrand \emph{before} applying
the theorem. A negative answer to a positive integrand, or an answer smaller
than an obvious lower bound, is the alarm --- and the sanity check that
catches it costs one glance. \chref{integrals}{6} is about doing this properly.}

\wstep{Step 6. The properties that save work.}
\[
  \int_{a}^{b}=-\int_{b}^{a},
  \qquad
  \int_{a}^{a}=0,
  \qquad
  \int_{a}^{b}=\int_{a}^{c}+\int_{c}^{b} .
\]
And by symmetry, over a symmetric interval $[-a,a]$: an \textbf{odd}
integrand gives $0$, and an \textbf{even} one gives $2\int_{0}^{a}$. So
$\int_{-2}^{2}x^{3}\mathrm{d}x=0$ without any computation --- provided the
integrand is genuinely defined throughout, which is exactly what failed in
Step~5.

\wstep{Step 7. Substitution in a definite integral, restated.}
Change the limits along with the variable, and the antiderivative never has
to be converted back. If the substitution is not monotone on the interval,
split the interval at the turning point first --- otherwise the two halves
cancel and the answer comes out too small.
\end{worked}

\section{Practice Problems}

Evaluate each definite integral; for the FTC-Part I problems,
differentiate the indicated function.

\begin{problems}
  \item $\displaystyle\int_{0}^{2} x^{2}\,\mathrm{d}x$
  \item $\displaystyle\int_{0}^{1} (3x^{2}+2x)\,\mathrm{d}x$
  \item $\displaystyle\int_{1}^{3} \frac{1}{x}\,\mathrm{d}x$
  \item $\displaystyle\int_{0}^{1} \mathrm{e}^{x}\,\mathrm{d}x$
  \item $\displaystyle\int_{0}^{2} (4-x^{2})\,\mathrm{d}x$
  \item $\displaystyle\int_{-1}^{1} (x^{3}+x)\,\mathrm{d}x$ \quad (symmetry?)
  \item $\displaystyle\int_{1}^{4} \sqrt{x}\,\mathrm{d}x$
  \item $\displaystyle\int_{0}^{1} \frac{x}{x^{2}+1}\,\mathrm{d}x$
  \item $\displaystyle\int_{0}^{\ln 2} \mathrm{e}^{2x}\,\mathrm{d}x$
  \item $\displaystyle\int_{1}^{\mathrm{e}} \frac{\ln x}{x}\,\mathrm{d}x$
  \item $\displaystyle\int_{0}^{1} x\,\mathrm{e}^{x^{2}}\,\mathrm{d}x$
  \item $\displaystyle\int_{1}^{2} \frac{1}{x^{2}}\,\mathrm{d}x$
  \item $\displaystyle\int_{0}^{3} (2x+1)^{2}\,\mathrm{d}x$
  \item $\displaystyle\int_{0}^{1} \frac{1}{2x+1}\,\mathrm{d}x$
  \item $\displaystyle\int_{0}^{4} \sqrt{2x+1}\,\mathrm{d}x$
  \item $\displaystyle\int_{0}^{1} x(1-x)^{4}\,\mathrm{d}x$
  \item Average value of $f(x)=x^{2}$ on $[0,3]$.
  \item Average value of $f(t)=\mathrm{e}^{-t}$ on $[0,1]$.
  \item Compute $\displaystyle\int_{2}^{0} (x+1)\,\mathrm{d}x$.
  \item $\displaystyle\int_{0}^{2} |x-1|\,\mathrm{d}x$ \quad (split at $x=1$)
  \item Let $A(x)=\displaystyle\int_{0}^{x} \sqrt{1+t^{2}}\,\mathrm{d}t$.
        Find $A'(x)$.
  \item Let $G(x)=\displaystyle\int_{0}^{x^{2}} \mathrm{e}^{-t^{2}}\,\mathrm{d}t$.
        Find $G'(x)$.
  \item $\displaystyle\int_{1}^{2} x\ln x\,\mathrm{d}x$
  \item $\displaystyle\int_{0}^{1} \frac{1}{(2x+1)^{2}}\,\mathrm{d}x$
  \item Total profit from marginal profit: $\Pi'(q)=16-0.3q$. Compute the
        total additional profit gained by raising output from $q=10$ to
        $q=30$.
\end{problems}

\section{Solutions}

\begin{solutions}
  \item $\displaystyle\Bigl[\frac{x^{3}}{3}\Bigr]_{0}^{2}=\frac{8}{3}$.

  \item $\displaystyle\bigl[x^{3}+x^{2}\bigr]_{0}^{1}=2$.

  \item $\displaystyle\bigl[\ln x\bigr]_{1}^{3}=\ln 3$.

  \item $\displaystyle\bigl[\mathrm{e}^{x}\bigr]_{0}^{1}=\mathrm{e}-1$.

  \item $\displaystyle\Bigl[4x-\frac{x^{3}}{3}\Bigr]_{0}^{2}
        =8-\frac{8}{3}=\frac{16}{3}$.

  \item Integrand is odd; interval symmetric about $0$. Therefore the
        integral is $0$. (Check:
        $\displaystyle\Bigl[\frac{x^{4}}{4}+\frac{x^{2}}{2}\Bigr]_{-1}^{1}
        =\bigl(\tfrac14+\tfrac12\bigr)-\bigl(\tfrac14+\tfrac12\bigr)=0$.)

  \item $\displaystyle\int_{1}^{4} x^{1/2}\,\mathrm{d}x
        =\Bigl[\frac{2}{3}x^{3/2}\Bigr]_{1}^{4}=\frac{2}{3}(8-1)=\frac{14}{3}$.

  \item Substitution $u=x^{2}+1$, $\mathrm{d}u=2x\,\mathrm{d}x$. Limits
        $u:1\to 2$.
        \[
          \int_{1}^{2} \frac{1}{2u}\,\mathrm{d}u=\frac{1}{2}\ln 2 .
        \]

  \item $\displaystyle\Bigl[\frac{1}{2}\mathrm{e}^{2x}\Bigr]_{0}^{\ln 2}
        =\frac{1}{2}\bigl(\mathrm{e}^{2\ln 2}-1\bigr)=\frac{1}{2}(4-1)=\frac{3}{2}$.

  \item $u=\ln x$, $\mathrm{d}u=\frac{1}{x}\,\mathrm{d}x$. Limits
        $u:0\to 1$.
        \[
          \int_{0}^{1} u\,\mathrm{d}u=\frac{1}{2} .
        \]

  \item $u=x^{2}$, $\mathrm{d}u=2x\,\mathrm{d}x$. Limits $u:0\to 1$.
        \[
          \int_{0}^{1} \frac{1}{2}\mathrm{e}^{u}\,\mathrm{d}u
            =\frac{1}{2}(\mathrm{e}-1) .
        \]

  \item $\displaystyle\Bigl[-\frac{1}{x}\Bigr]_{1}^{2}
        =-\frac{1}{2}+1=\frac{1}{2}$.

  \item $u=2x+1$, $\mathrm{d}u=2\,\mathrm{d}x$. Limits $u:1\to 7$.
        \[
          \int_{1}^{7} \frac{u^{2}}{2}\,\mathrm{d}u
            =\Bigl[\frac{u^{3}}{6}\Bigr]_{1}^{7}=\frac{343-1}{6}=57 .
        \]

  \item $\displaystyle\Bigl[\frac{1}{2}\ln(2x+1)\Bigr]_{0}^{1}
        =\frac{1}{2}\ln 3$.

  \item $u=2x+1$, limits $u:1\to 9$, $\mathrm{d}x=\frac{1}{2}\mathrm{d}u$.
        \[
          \int_{1}^{9} u^{1/2}\cdot\frac{1}{2}\,\mathrm{d}u
            =\Bigl[\frac{1}{3}u^{3/2}\Bigr]_{1}^{9}=\frac{27-1}{3}=\frac{26}{3} .
        \]

  \item $u=1-x$, so $x=1-u$, $\mathrm{d}x=-\mathrm{d}u$. Limits $u:1\to 0$.
        \[
          \int_{1}^{0} (1-u)u^{4}(-\mathrm{d}u)
            =\int_{0}^{1} (u^{4}-u^{5})\,\mathrm{d}u
            =\frac{1}{5}-\frac{1}{6}=\frac{1}{30} .
        \]

  \item $\displaystyle\bar f=\frac{1}{3}\int_{0}^{3} x^{2}\,\mathrm{d}x
        =\frac{1}{3}\cdot\frac{27}{3}=3$.

  \item $\displaystyle\bar f=\frac{1}{1}\int_{0}^{1} \mathrm{e}^{-t}\,\mathrm{d}t
        =\bigl[-\mathrm{e}^{-t}\bigr]_{0}^{1}=1-\mathrm{e}^{-1}\approx 0.632$.

  \item Reverse limits:
        $\displaystyle\int_{2}^{0} (x+1)\,\mathrm{d}x
        =-\int_{0}^{2} (x+1)\,\mathrm{d}x
        =-\Bigl[\frac{x^{2}}{2}+x\Bigr]_{0}^{2}=-(2+2)=-4$.

  \item Split at $1$:
        \[
          \int_{0}^{1} (1-x)\,\mathrm{d}x+\int_{1}^{2} (x-1)\,\mathrm{d}x
          =\Bigl[x-\frac{x^{2}}{2}\Bigr]_{0}^{1}
           +\Bigl[\frac{x^{2}}{2}-x\Bigr]_{1}^{2}
          =\frac{1}{2}+\bigl(0-(-\tfrac12)\bigr)=1 .
        \]

  \item By FTC I, $A'(x)=\sqrt{1+x^{2}}$.

  \item Chain rule + FTC I:
        $G'(x)=\mathrm{e}^{-(x^{2})^{2}}\cdot 2x=2x\,\mathrm{e}^{-x^{4}}$.

  \item By Section 3 problem 3, an antiderivative of $x\ln x$ is
        $\frac{x^{2}}{2}\ln x-\frac{x^{2}}{4}$. Evaluate:
        \[
          \Bigl[\frac{x^{2}}{2}\ln x-\frac{x^{2}}{4}\Bigr]_{1}^{2}
            =(2\ln 2-1)-\bigl(0-\tfrac14\bigr)=2\ln 2-\frac{3}{4} .
        \]

  \item $u=2x+1$, limits $u:1\to 3$.
        \[
          \int_{1}^{3} u^{-2}\cdot\frac{1}{2}\,\mathrm{d}u
            =\frac{1}{2}\bigl[-u^{-1}\bigr]_{1}^{3}
            =\frac{1}{2}\Bigl(-\frac{1}{3}+1\Bigr)=\frac{1}{3} .
        \]

  \item Total change in profit is
        $\displaystyle\int_{10}^{30} \Pi'(q)\,\mathrm{d}q$.
        \[
          \int_{10}^{30} (16-0.3q)\,\mathrm{d}q
            =\bigl[16q-0.15q^{2}\bigr]_{10}^{30}
            =(480-135)-(160-15)=345-145=200 .
        \]
\end{solutions}

\section{Additional Problems}

Evaluate each definite integral; for the FTC-Part I problems,
differentiate the indicated function.

\begin{problems}
  \item $\displaystyle\int_{0}^{3} x^{3}\,\mathrm{d}x$
  \item $\displaystyle\int_{1}^{2} (x^{2}+2)\,\mathrm{d}x$
  \item $\displaystyle\int_{1}^{4} \frac{1}{\sqrt{x}}\,\mathrm{d}x$
  \item $\displaystyle\int_{0}^{2} \mathrm{e}^{-x}\,\mathrm{d}x$
  \item $\displaystyle\int_{-2}^{2} (x^{2}-4)\,\mathrm{d}x$
  \item $\displaystyle\int_{-1}^{1} (x^{5}+x^{3})\,\mathrm{d}x$ \quad (symmetry?)
  \item $\displaystyle\int_{0}^{1} x\,\mathrm{e}^{x^{2}+1}\,\mathrm{d}x$
  \item $\displaystyle\int_{1}^{\mathrm{e}} \frac{(\ln x)^{2}}{x}\,\mathrm{d}x$
  \item $\displaystyle\int_{0}^{2} (3x^{2}-4x+1)\,\mathrm{d}x$
  \item $\displaystyle\int_{0}^{1} (4x+1)^{3}\,\mathrm{d}x$
  \item $\displaystyle\int_{0}^{1} \frac{x}{x^{2}+4}\,\mathrm{d}x$
  \item $\displaystyle\int_{0}^{\ln 3} \mathrm{e}^{-x}\,\mathrm{d}x$
  \item Average value of $f(x)=x^{3}$ on $[0,2]$.
  \item Average value of $f(x)=1/x$ on $[1,\mathrm{e}]$.
  \item $\displaystyle\int_{1}^{4} \frac{\sqrt{x}}{x}\,\mathrm{d}x$
  \item $\displaystyle\int_{3}^{0} (2x-1)\,\mathrm{d}x$
  \item $\displaystyle\int_{-1}^{3} |x|\,\mathrm{d}x$ \quad (split at $0$)
  \item Let $A(x)=\displaystyle\int_{1}^{x} \frac{1}{1+t^{4}}\,\mathrm{d}t$.
        Find $A'(x)$.
  \item Let $G(x)=\displaystyle\int_{0}^{x^{3}} \ln(1+t^{2})\,\mathrm{d}t$.
        Find $G'(x)$.
  \item $\displaystyle\int_{0}^{1} x^{2}\,\mathrm{e}^{x}\,\mathrm{d}x$
  \item $\displaystyle\int_{0}^{1} \frac{1}{(3x+2)^{2}}\,\mathrm{d}x$
  \item Marginal profit $\Pi'(q)=20-0.4q$. Additional profit from $q=20$
        to $q=50$.
  \item $\displaystyle\int_{0}^{2} (x+1)\,\mathrm{e}^{-x}\,\mathrm{d}x$
  \item $\displaystyle\int_{0}^{4} |x-2|\,\mathrm{d}x$
  \item Marginal cost $MC(q)=2+0.5q$. Total cost increment as $q$ goes
        from $0$ to $10$.
\end{problems}

\section{Additional Solutions}

\begin{solutions}
  \item $\displaystyle\Bigl[\frac{x^{4}}{4}\Bigr]_{0}^{3}=\frac{81}{4}$.

  \item $\displaystyle\Bigl[\frac{x^{3}}{3}+2x\Bigr]_{1}^{2}
        =\Bigl(\frac{8}{3}+4\Bigr)-\Bigl(\frac{1}{3}+2\Bigr)
        =\frac{7}{3}+2=\frac{13}{3}$.

  \item $\displaystyle\bigl[2\sqrt{x}\bigr]_{1}^{4}=4-2=2$.

  \item $\displaystyle\bigl[-\mathrm{e}^{-x}\bigr]_{0}^{2}
        =-\mathrm{e}^{-2}+1=1-\mathrm{e}^{-2}$.

  \item $\displaystyle\Bigl[\frac{x^{3}}{3}-4x\Bigr]_{-2}^{2}
        =\Bigl(\frac{8}{3}-8\Bigr)-\Bigl(-\frac{8}{3}+8\Bigr)
        =\frac{16}{3}-16=-\frac{32}{3}$.
        (Integrand $\le 0$ on $[-2,2]$.)

  \item Odd integrand on symmetric interval $\Rightarrow$ integral is $0$.

  \item $u=x^{2}+1$, $\mathrm{d}u=2x\,\mathrm{d}x$. Limits $u:1\to 2$.
        \[
          \int_{1}^{2} \frac{1}{2}\mathrm{e}^{u}\,\mathrm{d}u
            =\frac{1}{2}\bigl(\mathrm{e}^{2}-\mathrm{e}\bigr)
            =\frac{\mathrm{e}(\mathrm{e}-1)}{2} .
        \]

  \item $u=\ln x$, $\mathrm{d}u=\frac{1}{x}\,\mathrm{d}x$. Limits
        $u:0\to 1$.
        \[
          \int_{0}^{1} u^{2}\,\mathrm{d}u=\frac{1}{3} .
        \]

  \item $\displaystyle\bigl[x^{3}-2x^{2}+x\bigr]_{0}^{2}=8-8+2=2$.

  \item $u=4x+1$, $\mathrm{d}u=4\,\mathrm{d}x$. Limits $u:1\to 5$.
        \[
          \int_{1}^{5} \frac{u^{3}}{4}\,\mathrm{d}u
            =\Bigl[\frac{u^{4}}{16}\Bigr]_{1}^{5}=\frac{625-1}{16}=39 .
        \]

  \item $u=x^{2}+4$, $\mathrm{d}u=2x\,\mathrm{d}x$. Limits $u:4\to 5$.
        \[
          \int_{4}^{5} \frac{1}{2u}\,\mathrm{d}u=\frac{1}{2}\ln\frac{5}{4} .
        \]

  \item $\displaystyle\bigl[-\mathrm{e}^{-x}\bigr]_{0}^{\ln 3}
        =-\frac{1}{3}+1=\frac{2}{3}$.

  \item $\displaystyle\bar f=\frac{1}{2}\int_{0}^{2} x^{3}\,\mathrm{d}x
        =\frac{1}{2}\cdot 4=2$.

  \item $\displaystyle\bar f=\frac{1}{\mathrm{e}-1}\int_{1}^{\mathrm{e}}
        \frac{1}{x}\,\mathrm{d}x=\frac{1}{\mathrm{e}-1}$.

  \item $\displaystyle\frac{\sqrt{x}}{x}=x^{-1/2}$, so
        $\bigl[2\sqrt{x}\bigr]_{1}^{4}=2$.

  \item Reverse:
        $\displaystyle\int_{3}^{0} (2x-1)\,\mathrm{d}x
        =-\int_{0}^{3} (2x-1)\,\mathrm{d}x
        =-\bigl[x^{2}-x\bigr]_{0}^{3}=-(9-3)=-6$.

  \item Split at $0$:
        \[
          \int_{-1}^{0} (-x)\,\mathrm{d}x+\int_{0}^{3} x\,\mathrm{d}x
          =\frac{1}{2}+\frac{9}{2}=5 .
        \]

  \item By FTC I, $\displaystyle A'(x)=\frac{1}{1+x^{4}}$.

  \item Chain rule + FTC I: $G'(x)=\ln(1+x^{6})\cdot 3x^{2}$.

  \item Antiderivative of $x^{2}\,\mathrm{e}^{x}$ is
        $(x^{2}-2x+2)\,\mathrm{e}^{x}$ (Section 3 problem 5).
        \[
          \bigl[(x^{2}-2x+2)\,\mathrm{e}^{x}\bigr]_{0}^{1}
            =(1-2+2)\,\mathrm{e}-2=\mathrm{e}-2 .
        \]

  \item $u=3x+2$, $\mathrm{d}u=3\,\mathrm{d}x$. Limits $u:2\to 5$.
        \[
          \int_{2}^{5} \frac{1}{3u^{2}}\,\mathrm{d}u
            =\frac{1}{3}\Bigl[-\frac{1}{u}\Bigr]_{2}^{5}
            =\frac{1}{3}\Bigl(-\frac{1}{5}+\frac{1}{2}\Bigr)
            =\frac{1}{3}\cdot\frac{3}{10}=\frac{1}{10} .
        \]

  \item $\displaystyle\int_{20}^{50} (20-0.4q)\,\mathrm{d}q
        =\bigl[20q-0.2q^{2}\bigr]_{20}^{50}
        =(1000-500)-(400-80)=500-320=180$.

  \item Split:
        $\displaystyle\int_{0}^{2} x\,\mathrm{e}^{-x}\,\mathrm{d}x
        +\int_{0}^{2} \mathrm{e}^{-x}\,\mathrm{d}x$. By Section 3
        problem 2, $\int x\,\mathrm{e}^{-x}\,\mathrm{d}x=-(x+1)\,\mathrm{e}^{-x}$.
        \[
          \bigl[-(x+1)\,\mathrm{e}^{-x}\bigr]_{0}^{2}
           +\bigl[-\mathrm{e}^{-x}\bigr]_{0}^{2}
          =(-3\mathrm{e}^{-2}+1)+(-\mathrm{e}^{-2}+1)
          =2-4\mathrm{e}^{-2} .
        \]

  \item Split at $2$:
        \[
          \int_{0}^{2} (2-x)\,\mathrm{d}x+\int_{2}^{4} (x-2)\,\mathrm{d}x
          =\Bigl[2x-\frac{x^{2}}{2}\Bigr]_{0}^{2}
           +\Bigl[\frac{x^{2}}{2}-2x\Bigr]_{2}^{4}
          =2+2=4 .
        \]

  \item $\displaystyle\int_{0}^{10} (2+0.5q)\,\mathrm{d}q
        =\bigl[2q+0.25q^{2}\bigr]_{0}^{10}=20+25=45$.
\end{solutions}
